% FILE: 02_lineage_and_context.tex
% Project: livestream
% Thesis Role: real-time-process-broadcast-and-livestream-standards

\section{Lineage and Context}
\label{sec:livestream-lineage}

\subsection{2016 Language-Model Lesson}
\label{sec:livestream-lineage-2016}

The intellectual lineage of the livestream project begins with a 2016 observation
about local language-model imitation: a model that reproduces the surface form of
text does not thereby conserve the consequences of that text. A model can emit a
sentence that resembles a process description without any underlying process having
occurred. This observation crystallized a distinction that drives the entire
dissertation: \emph{prediction} (generating text that resembles process outcomes)
is categorically different from \emph{coordination} (emitting evidence that a
lawful process actually executed).

The livestream project is the most direct institutional expression of this lesson.
Its architecture refuses to treat the factory's outputs as self-certifying. Instead,
it requires that every significant factory event---every Git commit, every subagent
replay, every doctrine decision---be materialized as an OCEL 2.0 event and subjected
to conformance checking against a declared model. The model is not inferred from
the outputs; it is declared in advance (the Blue River Dam thesis Petri net), and
the conformance check asks whether the observed event sequence is consistent with
the declared model. This is precisely the structure of van der Aalst conformance
checking, applied not to a business process but to the factory that produces
process intelligence artifacts.

\subsection{Enterprise Process Gap}
\label{sec:livestream-lineage-enterprise-gap}

The enterprise process gap identified in the dissertation is the absence of
continuous, verifiable process evidence in software development organizations.
Software factories routinely produce artifacts---commits, builds, deployments,
test results---but these artifacts are not typically organized as a process-minable
event log. As a result, claims about the factory's behavior (e.g., ``we follow
a rigorous review process,'' ``every release is tested,'' ``no code ships without
conformance checking'') cannot be verified by an external auditor using process
mining methods. They are literary descriptions, not process evidence.

The livestream project addresses this gap directly. The \texttt{aalst\_broadcaster.js}
module transforms the Git commit history of the \texttt{process-intelligence}
repository into an OCEL 2.0 event stream: each commit becomes an event with
\texttt{ocel:activity}\,=\,``Git Commit'', \texttt{ocel:timestamp} derived from
the commit date (converted to UTC ISO 8601), and \texttt{ocel:vmap} fields carrying
the commit hash, author, and subject line. The author object (\texttt{ocel:omap})
links the event to an object in the OCEL 2.0 object graph. This transforms the
repository's history---previously legible only to humans reading \texttt{git log}
output---into a format that pm4py or wasm4pm can ingest for process discovery and
conformance checking.

\subsection{Relationship to the Chatman Equation}
\label{sec:livestream-lineage-chatman-equation}

\textbf{[AUTHOR THESIS]} The Chatman Equation $A = \mu(O^*, T, L)$ defines
admissible actuation as a function of the object model $O^*$, the declared
transition system $T$, and the observed event log $L$. The livestream project
instantiates this equation at the factory level: the factory's own actuation
(committing code, emitting receipts, authorizing downstream ggen manufacturing)
is governed by the MAPE-K autonomic controller in \texttt{autonomic.js}, which
evaluates the current event log against the declared Petri net model and withholds
authorization when conformance fitness falls below the lower control limit
(LCL\,=\,0.920).

This is not a metaphorical application of the Chatman Equation. The autonomic
controller implements the equation's structure directly: it monitors the event
stream ($L$), analyzes drift against the declared model ($T$), plans a remediation
action (rate limiting, dynamic route changes, S-component hot-swap), and executes
the plan through the MAPE-K tick function. The hot-swappable Petri net S-components
allow the declared model itself to be updated without stopping the factory, which
corresponds to updating $T$ in response to new information about the object model
$O^*$.

\subsection{Relationship to Receipts and Replay}
\label{sec:livestream-lineage-receipts-replay}

The receipt architecture of the process-intelligence foundry (documented in
\texttt{receipts/RECEIPT\_REGISTRY.md}) establishes that every major research
finding must be certified by a receipt: a named artifact that names what was
produced, by whom, under which witness, and what remains open. The livestream
project contributes to this architecture in two ways.

First, it provides the \emph{replay} surface: the \texttt{replay\_aggregator.js}
module consolidates all zoeapp AGI subagent replays from
\texttt{/Users/sac/zoeapp/replays/} into a single \texttt{master\_conversation.ocel}
file. This master file is the event log against which the conformance audit
(Fitness\,=\,0.9880, Precision\,=\,1.0000 against the Blue River Dam thesis model)
was computed. The conformance audit result is itself a receipt: it certifies that
the AGI conversation trace is consistent with the declared thesis process model.

Second, the AuditLedger in the browser visualizer implements cryptographic receipt
chaining at the event level: each ledger block carries a SHA-256 hash that chains
to the previous block, so that any tampering with a historical event is immediately
detectable. This extends the receipt architecture from the document level (where
receipts are YAML or JSON files in the \texttt{receipts/} directory) to the event
level (where every process event is individually ledgered and chained).
